\section{一般框架}\label{sect_formalism}%
%

%
\subsection{部分子准分布}\label{subsect_definitions}%
%
我们从如下核子矩阵元出发
\begin{align}
\label{def:ME}
 &\langle N(p)| \bar q (zv) \gamma_\alpha[zv,0] q(0) |N(p)\rangle  =
\phantom{\frac{2v_\alpha}{v^2} (pv) \mathcal{I}^\parallel (v^2z^2\!,pvz,\mu)}
\notag\\
&\hspace*{2.0cm} = \frac{2v_\alpha}{v^2} (pv)\, \mathcal{I}^\parallel (v^2z^2\!,pvz,\mu)
\notag\\
&\hspace*{2.0cm}\quad+ 2\left(p_\alpha- \frac{v_\alpha}{v^2} (pv) \right) \mathcal{I}^\perp (v^2z^2\!,pvz,\mu)\,,
\end{align}
其中
\begin{align}
\label{def:Wilson}
 [zv, 0] &= \text{Pexp}\Big[i g \int_0^z du\,  v^\mu A_\mu(uv)\Big]\,,
\end{align}
$p^\mu$ 是核子动量，$p^2 = m^2$，$v^\mu$ 是给定的四维矢量，$z$ 为实数。所有记号均对应 Minkowski 空间。下文中我们保持四矢 $v_\mu$ 的归一化任意，但记住在格点计算中 $v^2 <0$。我们省略味道指标，并默认在以下讨论中考虑矩阵元的味道非单态组合。我们还忽略夸克质量。

式 \eqref{def:ME} 中的算符乘积存在紫外（UV）发散，必须重正化。矩阵元 $\mathcal{I}$ 的宗量 $\mu$ 指重正化尺度。重正化尺度依赖可以通过过渡到有效理论来研究~\cite{Craigie:1980qs,Dorn:1986dt,Ji:2017oey}，此时 Wilson 线被一个辅助场的传播子所替代。对于类时间隔，所得理论是重夸克有效理论（HQET），所关心的重正化因子就是重—轻夸克流的重正化因子的平方，见例如~\cite{Chetyrkin:2003vi}。据我们所知，此类重正化因子在类空间隔下尚无单圈以上的计算，不过在该精度下与类时情形没有差别。


式~\eqref{def:ME} 中的``纵向''与``横向''（相对 $v_\mu$ 而言）不变函数对应于格点计算中所采用的矩阵元的特定投影：
\begin{align}\label{invfuncdef}
 \langle N(p)| \bar q (zv) \slashed{v} [zv,0]q(0) |N(p)\rangle  & = 2 (pv) \mathcal{I}^\parallel (v^2z^2,pvz,\mu)\,,
\notag\\
  \langle N(p)| \bar q (zv) \slashed{\epsilon}[zv,0] q(0) |N(p)\rangle &= 2 (p\epsilon) \mathcal{I}^\perp (v^2z^2,pvz,\mu)\,,
\end{align}
其中 $(\epsilon\cdot v) =0$。式~(\ref{invfuncdef}) 是构造准分布的出发点。

不变函数 $\mathcal{I}^{\parallel}$、$\mathcal{I}^{\perp}$ \eqref{def:ME} 在树图层面相同。假设幂计数
\begin{align}
\label{power-counting}
  z = \mathcal{O}(\eta)\,, \qquad p = \mathcal{O}(\eta^{-1})\,,\qquad \eta \to 0\,,
\end{align}
则它们可以用位置空间的夸克 PDF~\cite{Ioffe:1969kf,Braun:1994jq,Hoyer:1996nr} 写成：
\begin{align}
  \mathcal{I}^{\parallel(\perp)}(z^2v^2,pvz, \mu\sim 1/|vz|) &=  I(pvz, \mu_F \sim 1/|vz|)
%\int_{-1}^1 dx \, e^{ixz\,pv} q(x,\mu_F \sim 1/|z|)
\notag\\&\quad
 +\mathcal{O}(\alpha_s, \eta^2)\,,
\end{align}
其中 $\mu_F$ 是因子化尺度，且
\begin{align}
\label{ITD}
   I(pvz,\mu_F) &=  \int_{-1}^1 dx \, e^{ixz\,pv} q(x,\mu_F)\,.
\end{align}
这里 $q(x,\mu_F)$ 在 $x>0$ 时是夸克 PDF，在 $x<0$ 时是反夸克 PDF。位置空间 PDF $I$ 被称为 Ioffe 时间分布（ITD）。为了把 ITD $I$ 与函数 $\mathcal{I}$ 区分开，按照文献 \cite{Ji:2013dva} 引入的术语，我们把函数 $\mathcal{I}^{\parallel}$、$\mathcal{I}^{\perp}$ 称为 Ioffe 时间\emph{准}分布（qITD）。

qPDF $\mathcal{Q}^{\parallel},\mathcal{Q}^{\perp}$~\cite{Ji:2013dva} 和 pPDF $\mathcal{P}$~\cite{Radyushkin:2017cyf}
由 qITD 通过 Fourier 变换定义：
\begin{align}
\label{def:qPDF+pPDF}
\mathcal{Q}^{\parallel(\perp)}(x,p,\mu) &= |(pv)|\! \int\limits_{-\infty}^\infty \frac{d z }{2\pi} e^{-ix z (pv)} \, \mathcal{I}^{\parallel(\perp)}(z^2v^2, pvz,\mu )\,,
\notag\\
\mathcal{P}(x,z,\mu) &= |z| \int\limits_{-\infty}^\infty \frac{d (pv)}{2\pi} e^{-ix z (pv) } \,  \mathcal{I}^{\perp}(z^2v^2,pvz,\mu)\,.
\end{align}
下文中我们默认在讨论 qPDF 时 $(pv)>0$，讨论 pPDF 时 $z>0$，并为简洁省去绝对值符号。

在具有显式正规化尺度的重正化方案中，式 \eqref{def:ME} 中的 Wilson 线还受到额外的线性 UV 发散，必须予以去除。在维数正规化中，这一线性 UV 发散表现为微扰级数高阶项的阶乘增长~\cite{Beneke:1994sw}。线性 UV 发散可以通过与 Wilson 线相联系的质量重正化去除~\cite{Dotsenko:1979wb,Craigie:1980qs,Dorn:1986dt}，也可以通过正规化无关的重正化方案去除~\cite{Martinelli:1995vj,Stewart:2017tvs,Chen:2017mzz}。鉴于准 PDF 算符的乘法可重正性~\cite{Ji:2017oey,Ishikawa:2017faj,Green:2017xeu}，还可以通过构造涉及同一算符的矩阵元的适当比值将其消除，例如按零质子动量处的同一矩阵元归一化~\cite{Orginos:2017kos}，
\begin{align}
\label{n-qITD}
     \mathbf{I}(v^2z^2, pvz) =  \mathcal{I}(v^2z^2, pvz, \mu) / \mathcal{I}(v^2z^2, 0, \mu)\,,
\end{align}
或者改为按真空期望值归一化：
\begin{align}
\label{hat-qITD}
     \widehat{\mathbf{I}}(v^2z^2, pvz) =  \mathcal{I}(v^2z^2, pvz, \mu) / \mathcal{N}(v^2z^2,\mu)\,,
\end{align}
其中
\begin{align}
  \mathcal{N}(v^2z^2,\mu) =  \left(\frac{2 i N_c}{\pi^2 z^3 v^2}\right)^{-1} \langle 0| \bar q(zv) \slashed{v} [zv,0] q(0) |0\rangle\,.
\end{align}
式~(\ref{n-qITD}) 与 (\ref{hat-qITD}) 是本文的主要关注点。通过作比值，尺度依赖（包括通常的对数重正化）被消去，于是可以定义与尺度无关的 qPDF/pPDF：
\begin{align}
\mathbf{Q}(x,p) &= |(pv)| \int_{-\infty}^\infty \frac{d z }{2\pi} e^{-ix z (pv)} \, \mathbf{I}(z, pv)\,,
\notag\\
\mathbf{P}(x,z) &= |z| \int_{-\infty}^\infty \frac{d (pv)}{2\pi} e^{-ix z (pv) } \,  \mathbf{I}(z,pv)\,,
\end{align}
对 $\widehat{\mathbf{Q}}(x,p)$ 与 $\widehat{\mathbf{P}}(x,z)$ 亦类似。
在本文语境下，这两种选择的区别在于：对真空关联函数的归一化不影响作为本文研究对象的领头 $\mathcal{O}(v^2z^2)$ 幂修正（因为不存在相应的规范不变算符），而我们将看到，按零动量取值的归一化则有可观的影响。
%Alternatively, the linear UV divergence can be tamed by  the nonperturbative subtractions, see e.g.~\cite{Martinelli:1995vj}.

%
\subsection{光线 OPE 与靶质量修正}\label{subsect_OPE}%
%

QCD 振幅位置空间共线因子化的一般方法由光线 OPE（light-ray operator product expansion）提供~\cite{Anikin:1978tj,Anikin:1979kq,Balitsky:1987bk,Mueller:1998fv,Balitsky:1990ck,Geyer:1999uq}。为便于说明，考虑``纵向''投影。针对当前情形（前向矩阵元），我们写出
\begin{align}
\label{light-ray-OPE}
 \bar q(zv) \slashed{v} [zv,0] q(0) &= \int\limits_0^1\!d\alpha \, H^\parallel(z,\alpha,\mu_F)
\notag\\&\quad \times \Pi_{\rm l.t.}^{\mu_F}[\bar q (\alpha z v) \slashed{z} q(0)] + \ldots,
\end{align}
其中重正化因子已包含在系数函数 $H^\parallel = \delta(1-\alpha) + \mathcal{O}(\alpha_s)$ 中，并且取重正化尺度 $\mu$ 等于因子化尺度 $\mu_F$。
最后，$\Pi_{\rm l.t.}^{\mu_F}[\ldots]$ 是下文定义的领头扭度投影算符，省略号代表高扭度贡献。

非定域夸克—反夸克算符的领头扭度投影定义为\emph{重正化的}定域领头扭度算符（无迹且对所有指标对称化）的生成函数：
\begin{align}
\Pi_{\rm l.t.}^{\mu_F}[\bar q (z v) \slashed{v} q(0)]
&= \sum_{n=1}^\infty \frac{z^{n-1}}{(n-1)!} v^{\mu_1}\ldots v^{\mu_n} O_{\mu_1\ldots\mu_n}^{n}(0)\,,
\label{eq:LTprojector}
\end{align}
其中
\begin{align}
  O_{\mu_1\ldots\mu_n}^{n}(z) &=
\bar q(0)\gamma_{(\mu_1}\!\!\stackrel{\leftarrow}{D}_{\mu_2}\ldots \!\stackrel{\leftarrow}{D}_{\mu_{n})} q(0)\,.
\end{align}
%
这里我们用圆括号括起相关 Lorentz 指标来表示减除迹与对称化，
例如 $O_{(\mu\nu)} = \frac12 (O_{\mu\nu}+ O_{\nu\mu}) - \frac14 g_{\mu\nu} O_{\lambda}^{~\lambda}$。

光线 OPE 与通常基于定域算符的短距离 Wilson 展开的区别在于采用了不同的幂计数。在短距离展开中，假定夸克间距很小，$|z| \sim \eta \Lambda^{-1}_{\rm QCD}$ 且 $\eta\to 0$，而在此极限下算符矩阵元为单位量级，$\langle O_{\mu_1\ldots\mu_n}^{n} \rangle \sim \Lambda^{n}_{\rm QCD}$。此时只有有限多项贡献到式~\eqref{eq:LTprojector} 右端，其余各项以及高扭度算符必须从 $\mathcal{O}(\eta^2)$ 开始逐阶加入 OPE 的高阶项。换言之，相关的展开参数是算符的量纲。光线 OPE 则假定领头扭度算符按 $\langle O_{\mu_1\ldots\mu_n}^{n} \rangle \sim \eta^{-n} \Lambda^{n}_{\rm QCD}$ 标度，使得 $ z^n v^{\mu_1}\ldots v^{\mu_n} \langle O_{\mu_1\ldots\mu_n}^{n} \rangle  = \mathcal{O}(1)$。此时式~\eqref{eq:LTprojector} 中的级数必须全阶求和，从而显现出其非定域的``光线''结构。对领头扭度算符的一般强子矩阵元有
\begin{align}
 \langle p | O_{\mu_1\ldots\mu_n}^{n} |p\rangle \sim p_{(\mu_1}\ldots p_{\mu_n)} \langle\!\langle O^{n} \rangle\!\rangle\,,
\end{align}
其中约化矩阵元 $\langle\!\langle O^{n} \rangle\!\rangle = \mathcal{O}(1)$。因此，若强子具有大动量 $|pv| = \mathcal{O}(\eta^{-1})$ 从而 $ z\,pv =  \mathcal{O}(1)$，光线 OPE 就是适当的近似。同一维度的高扭度算符自旋较小（按定义），因而其矩阵元受幂次压制\footnote{严格地说，大动量的幂展开对应于所谓的共线扭度分类，见例如~\cite{Geyer:1999uq,Geyer:2000ig}。}。

注意上述幂计数在 Minkowski 空间和 Euclidean 空间都适用。在 Minkowski 空间中，可以变换到一个新的参考系，使动量的所有分量都是 $\Lambda_{\rm QCD}$ 量级，同时夸克间距几乎类光：$|v_\mu| = \mathcal{O}(1)$ 但
$v^2 = \mathcal{O}(\eta^2) \to 0$。

就计算而言，光线 OPE 提供了一个处理领头扭度投影算符 \eqref{eq:LTprojector} 而避开定域展开的便利框架。光线算符可以看作夸克间距的解析算符函数（所有短距离奇点与光锥奇点均已被减除）。它们满足 Laplace 方程~\cite{Balitsky:1987bk}
\begin{align}
 \frac{\partial}{\partial v^\mu} \frac{\partial}{\partial v_\mu}  \Pi_{\rm l.t.}^{\mu_F}[\bar q (z v ) \slashed{v} q(0)] =0\,
\end{align}
边界条件为当 $v^2 \to 0$ 时 $\Pi_{\rm l.t.}^{\mu_F}\to 1$。
投影算符 $\Pi_{\rm l.t.}^{\mu_F}$ 的显式表达式可在文献~\cite{Balitsky:1987bk,Geyer:1999uq,Geyer:2000ig,Braun:2011dg} 中找到。
光线 OPE 结合背景场方法是标准技术，例如
在光锥求和规则~\cite{Balitsky:1989ry}中被用于计算高扭度贡献，
并用于推导离前向部分子分布的演化方程~\cite{Belitsky:1999hf,Belitsky:1998gc}。
一般离前向运动学中味道单态光线算符演化的重正化群核已知到三圈精度~\cite{Braun:2017cih}。

领头扭度投影算符 \eqref{eq:LTprojector} 的核子矩阵元定义了领头扭度
夸克 PDF：
\begin{eqnarray}
\lefteqn{ \langle N(p)| \Pi_{\rm l.t.}^{\mu_F}[\bar q (zv) \slashed{v} q(0)] |N(p)\rangle}
\notag\\
&=& 2 \int_{-1}^1\!dx\, \Pi_{\rm l.t.}[(p\cdot v) e^{ i x z (pv)}] q(x,\mu_F)\,,
\end{eqnarray}
其中~\cite{Balitsky:1990ck}
\begin{align}
\label{Pi-exp}
\Pi_{\rm l.t.}[(pv) e^{ i x z(pv)}]
&= (pv)\Big[1 - \frac{1}{4} \frac{m^2 v^2}{(pv)^2} z\frac{d}{dz} \Big] e^{ ixz(pv)}
\nonumber\\&\quad
+ \mathcal{O}(m^4)\,.
\end{align}
方括号中的第二项是核子质量的领头修正，它是深度非弹性散射（DIS）中 Nachtmann 靶质量修正~\cite{Nachtmann:1973mr} 在位置空间的对应物。领头扭度投影指数函数关于核子质量的全阶表达式见文献~\cite{Balitsky:1990ck}。

对式 \eqref{light-ray-OPE} 的算符关系取核子矩阵元，便得到``纵向'' qITD 的因子化定理：
\begin{align}
\mathcal{I}^\parallel   = &
 \int\limits_0^1 d\alpha \, H^\parallel(z,\alpha,\mu_F)
  \Big[1 - \frac{1}{4} \frac{m^2 v^2}{(pv)^2} z\frac{d}{dz} +\ldots \Big]\\ & \nonumber
\times I(\alpha z pv , \mu_F)
+\ldots\,,
\end{align}
其中 $I(\alpha z pv,\mu_F)$ 是 ITD \eqref{ITD}，省略号代表更高阶的靶质量修正以及``本征''高扭度修正。

作 Fourier 变换 \eqref{def:qPDF+pPDF}，在领头扭度精度下得到``纵向'' qPDF 的因子化定理：
\begin{align}
\label{qPDF2}
 \mathcal{Q}^\parallel(x,p) &= \int_{-1}^1 \!\frac{dy}{|y|} C^\parallel_{\Qd}(\tfrac{x}{y},xp,\mu_F)
\Big[1 + \frac{1}{4} \frac{m^2 v^2}{(pv)^2} \frac{d}{dy} y +\ldots \Big]
\notag\\&\quad
\times  q(y,\mu_F)\,,
\end{align}
其中系数函数为
\begin{align}
  C^\parallel_{\Qd}(\tfrac{x}{y},xp,\mu_F) &= (pv)|y|\!\int\limits_{-\infty}^{\infty}\!\frac{dz}{2\pi}\! \int\limits_0^1\!\! d\alpha\,
e^{i(pv)z(x-\alpha y)}
\notag\\&\quad \times H^\parallel(z,\alpha,\mu_F)\,.
\end{align}
特别地，在领头阶
\begin{align}
\label{target-parallel}
 \mathcal{Q}^\parallel(x,p) &= q(x) + \frac{1}{4} \frac{m^2 v^2}{(pv)^2} \big[x q'(x)  + q(x)\big]
+ \mathcal{O}\left({m^4}/{p^4}\right),
\end{align}
其中 $q'(x) = (d/dx) q(x)$。注意作用在夸克 PDF 上的导数会把幂次行为 $q(x) \stackrel{x \to 1}{\sim} (1-x)^p$ 降低一个幂次，因此在大 Bjorken $x$ 处，质量修正相对于夸克分布本身被有效放大了因子 $1/(1-x)$。DIS 中靶质量修正在 $x\to 1$ 处的类似增强是众所周知的~\cite{Nachtmann:1973mr}。

``横向'' qPDF 的靶质量修正可以用类似方式计算：
从带开放 Lorentz 指标的非定域算符 \eqref{def:ME} 出发，
利用算符恒等式~\cite{Balitsky:1987bk}
\begin{align}
\Pi_{\rm l.t.}^{\mu_F}[\bar q (z v) \gamma_\alpha q(0)]
=  \int\limits_0^\infty \!dt\,\frac{\partial}{\partial v^\alpha}
\Pi_{\rm l.t.}^{\mu_F}[\bar q ( t z v) \slashed{v} q(0)]\,,
\end{align}
其中夸克之间的 Wilson 线不言而喻。在树图层得到
\begin{align}
\label{target-perp}
 \mathcal{Q}^\perp(x,p)  &=  q(x) + \frac{1}{4} \frac{ m^2 v^2 }{(pv)^2} \Big[ x q'(x) + 3 q(x) \Big]
\notag\\&\quad
-  \frac{1}{2} \frac{m^2 v^2}{(pv)^2} \theta(|x|<1)\!\! \int_{|x|}^1\! \frac{dy}{y} q(x/y)
+ \mathcal{O}\left({m^4}/{p^4}\right),
\end{align}
而 pPDF 的靶质量修正为
\begin{align}
\label{target-pPDF}
 \mathcal{P}(x,z) &=  q(x) +  \frac14 z^2 v^2 m^2 x^2 \theta(|x|<1) \int_{|x|}^1 \frac{dy}{y} q(x/y)
\notag\\&\quad
+ \mathcal{O}\left({m^4}/{p^4}\right)\,.
\end{align}
将式~(\ref{target-parallel}) 和 (\ref{target-perp}) 展开到 $\mathcal O(m^2v^2/(pv)^2)$ 阶，其结果与文献~\cite{Chen:2016utp} 中推导的一致。有趣的是，pPDF 的靶质量修正在 $x\to 1$ 处以 $\mathcal{O}(1-x)$ 压低，而不像 qPDF 的修正 (\ref{target-parallel}) 及 DIS 结构函数那样被放大 $\mathcal{O}(1/(1-x))$。

\subsection{Borel 变换与 renormalon}

式 \eqref{light-ray-OPE} 中的系数函数 $H^\parallel$ 以及 $\MSbar$ 方案中类似的系数函数
$H^\perp$ 有如下微扰展开：
\begin{align}
 H = \delta(1-\alpha) + \sum_{k=0}^\infty h_k a_s^{k+1}\,,\qquad a_s = \frac{\alpha_s(\mu)}{4\pi}\,,
\end{align}
系数 $h_k\sim k!$ 阶乘增长。

处理这类级数的便捷方式是考察 Borel 变换
\begin{align}
  B[H] (w) = \sum_{k=0}^\infty \frac{h_k}{k!}\left(\frac{w}{\beta_0}\right)^k
\end{align}
其中插入 $\beta_0 = 11/3 N_C -2/3 n_f$ 的幂是为了后文的方便。Borel 像可用作固定阶系数的生成函数
\begin{align}
 h_k = \beta_0^k \left(\frac{d}{dw}\right)^k  B[H] (w)\big|_{w=0}\,.
\end{align}
此外，级数之和可以表示为对 Borel 参数 $w$ 正值区间的积分
\begin{align}
\label{BorelIntegral}
H = \delta(1-\alpha)  + \frac{1}{\beta_0} \int_0^\infty \!dw\, e^{-w/(\beta_0 a_s)} B[H](w)\,.
\end{align}
如上所述，该积分并无良好定义，因为 Borel 变换在积分路径上一般存在奇点，即所谓的（红外）renormalon。可以对积分采取变形围道（从实轴上方或下方绕过）的定义，或取主值。这些定义都是任意的，它们的差在耦合常数下是指数小量，必须视为微扰论的内在不确定性，需要通过添加幂次压低的非微扰修正来消除。另一个潜在问题涉及 Borel 积分在
$w\to\infty$ 处的收敛性。由于所研究的量只依赖于单一硬尺度 $1/|z v|$，量纲分析要求 Borel
变换可以写成 $(\mu^2 z^2 v^2 )^w$ 乘以 Borel 参数和无量纲运动学变量的函数 $F(w)$。
将 $(\mu^2 z^2 v^2)^w$ 与 $e^{-w/(\beta_0 a_s(\mu))}$ 合并可见，（主值）积分收敛的条件是
夸克间距 $|zv|$ 远小于 $1/\Lambda_{\rm QCD}$，且 $F(w)$ 在 $w\to\infty$ 处不呈指数增长——所有已知例子都满足后者。

自然，完整的全阶计算无法完成。我们转而采用近似~\cite{Beneke:1998ui,Beneke:2000kc}：
只考虑由单圈图中跑动耦合效应生成的微扰级数，即在胶子虚度能标处取 QCD 耦合。
这类贡献可以通过计算单圈图中插入 $k$ 个费米子圈的图
并把 $- \frac23 n_f \mapsto   \beta_0 = \frac{11}{3} N_c - \frac23 n_f$ 来追踪，见附录 A。
另一种等价技术 \cite{Ball:1995ni} 基于带有效胶子质量的单圈图计算。

%
\begin{figure*}[tbp]
\centering
 \includegraphics[width=0.95\textwidth]{BVZplots/bubblechain}
  %~\includegraphics[width=6.0cm]{BMJplots/fit-t4-qPDF}
\caption{气泡链对系数函数的贡献。Wilson 线因子用双点线表示。}
\label{fig:bubblechain}
\end{figure*}
%

Borel 变换的奇点结构可以不必显式计算气泡链而被单独提取出来。做法是把圈图中的跑动耦合常数替换为其有效形式：
\begin{align}
  \beta_0 a_s(-k^2) =  \int_0^\infty \!dw\, e^{\frac53 w} \left(\frac{\Lambda^2}{-k^2}\right)^w ,
\end{align}
其中 $\Lambda \equiv \Lambda_{\rm QCD}^{\MSbar}$，因子 $e^{\frac53 w}$ 代表 $\MSbar$ 方案。这样的替换导致胶子传播子变为
\begin{align}\label{modgluonprop}
  \frac{1}{-k^2-i\epsilon} ~\mapsto~ \frac{(\Lambda^2)^w}{(-k^2-i\epsilon)^{1+w}}\,.
\end{align}
在单圈层面，$w$ 平面上的极点结构重现了 Borel 平面的极点结构。在本工作中我们对结果做了两种方法的交叉检验。


%
%%%%%%%%%%%%%%%%%%%%%%%%%%%%%%%%%%%%%%%%%%%%%%%%%%%%%
